\section{Hypercontractivity --- Random Variables --- Sharp --- Two Point Law}
%%% generated by the blueprint pipeline from TCSlib.BooleanAnalysis.Hypercontractivity.RandomVariables.Sharp.TwoPointLaw
\subsection*{Hypercontractivity --- Random Variables --- Sharp --- Two Point Law: Introduction}
The sharp scalar bounds give affine hypercontractivity for the centered two-point distribution. The conjugate inequality also forces the radius to be no larger than the stated sharp value.

\subsection*{Hypercontractivity --- Random Variables --- Sharp --- Two Point Law: Results}

\begin{theorem}[Sharp conjugate contraction for a two point law]
\lean{BooleanAnalysis.Hypercontractivity.SharpDiscrete.two_point_conjugate_sharp}
\leanok
\uses{BooleanAnalysis.Hypercontractivity.FiniteLaw.ae_two_point_of_mass_sum, BooleanAnalysis.Hypercontractivity.FiniteLaw.two_point_affine_norm_sq, BooleanAnalysis.Hypercontractivity.HasAtomLowerBound, BooleanAnalysis.Hypercontractivity.HasAtomLowerBound.memLp, BooleanAnalysis.Hypercontractivity.IsHypercontractive, BooleanAnalysis.Hypercontractivity.SharpDiscrete.radius_pos_lt_one, BooleanAnalysis.Hypercontractivity.SharpDiscrete.weighted_biased_two_point_le, BooleanAnalysis.Hypercontractivity.rvLpNorm, BooleanAnalysis.Hypercontractivity.sharpDiscreteRadius}
\difficulty{5}
Let $(\Omega,\mu)$ be a probability space, let $X:\Omega\to\bbr$ be measurable, and
suppose $q>2$ and $0<\lambda<1/2$. Assume $\mu(X=1-\lambda)=\lambda$ and
$\mu(X=-\lambda)=1-\lambda$. Then $X$ is affine hypercontractive from exponent $q/(q-1)$
to exponent $2$ with parameter $\rho(q,\lambda)$, the sharp discrete hypercontractive
radius; in particular, this means $X\in L^2(\mu)$, $0\le\rho(q,\lambda)<1$, and
$\|a+\rho(q,\lambda)bX\|_2\le\|a+bX\|_{q/(q-1)}$ for all $a,b\in\bbr$.
\end{theorem}

\begin{theorem}[Upper bound from a two point law]
\lean{BooleanAnalysis.Hypercontractivity.SharpDiscrete.two_point_conjugate_radius_le}
\leanok
\uses{BooleanAnalysis.Hypercontractivity.FiniteLaw.two_point_affine_norm_sq, BooleanAnalysis.Hypercontractivity.IsHypercontractive, BooleanAnalysis.Hypercontractivity.SharpDiscrete.weighted_biased_radius_le, BooleanAnalysis.Hypercontractivity.rvLpNorm, BooleanAnalysis.Hypercontractivity.sharpDiscreteRadius}
\difficulty{5}
Let $(\Omega,\mu)$ be a probability space, let $X:\Omega\to\bbr$ be measurable, and
suppose $q>2$ and $0<\lambda<1/2$. Assume $\mu(X=1-\lambda)=\lambda$ and
$\mu(X=-\lambda)=1-\lambda$. If $X$ is affine hypercontractive from exponent $q/(q-1)$
to exponent $2$ with parameter $\rho$—that is, $0\le\rho<1$, $X\in L^2(\mu)$, and
$\|a+\rho bX\|_2\le\|a+bX\|_{q/(q-1)}$ for all $a,b\in\bbr$—then
$\rho\le\rho(q,\lambda)$, where $\rho(q,\lambda)$ is the sharp discrete hypercontractive
radius.
\end{theorem}

\begin{theorem}[Sharp forward contraction for a two point law]
\lean{BooleanAnalysis.Hypercontractivity.SharpDiscrete.two_point_forward_sharp}
\leanok
\uses{BooleanAnalysis.Hypercontractivity.FiniteLaw.ae_two_point_of_mass_sum, BooleanAnalysis.Hypercontractivity.FiniteLaw.two_point_affine_norm_sq, BooleanAnalysis.Hypercontractivity.HasAtomLowerBound, BooleanAnalysis.Hypercontractivity.HasAtomLowerBound.memLp, BooleanAnalysis.Hypercontractivity.IsHypercontractive, BooleanAnalysis.Hypercontractivity.SharpDiscrete.radius_pos_lt_one, BooleanAnalysis.Hypercontractivity.SharpDiscrete.weighted_biased_forward_le, BooleanAnalysis.Hypercontractivity.rvLpNorm, BooleanAnalysis.Hypercontractivity.sharpDiscreteRadius}
\difficulty{5}
Let $(\Omega,\mu)$ be a probability space, let $X:\Omega\to\bbr$ be measurable, and
suppose $q>2$ and $0<\lambda<1/2$. Assume $\mu(X=1-\lambda)=\lambda$ and
$\mu(X=-\lambda)=1-\lambda$. Then $X$ is affine hypercontractive from exponent $2$ to
exponent $q$ with parameter $\rho(q,\lambda)$, the sharp discrete hypercontractive
radius; that is, $X\in L^q(\mu)$, $0\le\rho(q,\lambda)<1$, and
$\|a+\rho(q,\lambda)bX\|_q\le\|a+bX\|_2$ for all $a,b\in\bbr$.
\end{theorem}
